\section{Evaluation}\label{sec:eval}

We evaluate \thunderbolt along three axes: (1) per-stage microbenchmarks of a single off-chain transfer (\autoref{sec:eval-micro}), (2) a system-level comparison against the Lightning Network (\autoref{sec:eval-lightning}), and (3) a comparison of our Sigma-protocol construction against general-purpose SNARK backends (\autoref{sec:eval-snark}).

\paragraph{Experimental setup.}
All benchmarks were run on a single machine with an Apple~M4 processor (10-core, 4 performance + 6 efficiency), 16\,GB unified memory, macOS~15.3, compiled with \texttt{rustc}~1.87-nightly in release mode with LTO enabled. The \thunderbolt implementation uses the \texttt{rust-paillier} crate for Paillier operations with a 3072-bit modulus ($|N| = 3072$, so all Paillier arithmetic operates over 6144-bit integers in $\ZZ_{N^2}$) and the \texttt{k256} crate for secp256k1 elliptic-curve arithmetic. Each experiment reports the average, minimum, and maximum over 100 iterations, with a 3-iteration warmup discarded before measurement.

\subsection{Per-Stage Microbenchmark}\label{sec:eval-micro}

A single off-chain transfer executes nine stages in sequence: invoice setup, invoice prove, invoice verify, transfer setup, transfer prove, anchor verify, transfer verify, Bob decrypt, and committee update. \autoref{tab:stages} reports the measured wall-clock time for each stage.

\begin{table}[t]
\centering
\small
\begin{tabular}{@{}lrrr@{}}
\toprule
\textbf{Stage} & \textbf{Avg (ms)} & \textbf{Min (ms)} & \textbf{Max (ms)} \\
\midrule
Invoice setup       & 118.1 & 114.4 & 143.2 \\
Invoice prove       & 305.3 & 294.5 & 351.3 \\
Invoice verify      & 194.1 & 187.0 & 232.9 \\
Transfer setup      &  63.3 &  60.9 &  68.1 \\
Transfer prove      & 159.3 & 154.3 & 178.6 \\
Anchor verify       &$<\!0.01$& $<\!0.01$ & $<\!0.01$ \\
Transfer verify     & 109.0 & 105.5 & 117.5 \\
Bob decrypt         &  30.2 &  28.6 &  34.0 \\
Committee update    &  30.2 &  29.1 &  34.8 \\
\midrule
\textbf{Total (off-chain)} & \textbf{1009.6} & \textbf{978.6} & \textbf{1089.9} \\
\bottomrule
\end{tabular}
\caption{Per-stage wall-clock time for a single \thunderbolt off-chain transfer (100 iterations, 3072-bit Paillier modulus, secp256k1).}
\label{tab:stages}
\end{table}

\paragraph{Cost breakdown.} The total off-chain transfer completes in \textbf{$\approx\!1.01$\,s} on average. \autoref{tab:breakdown} groups the nine stages by their functional role: proving, verification, setup, and auxiliary operations.

\begin{table}[t]
\centering
\small
\begin{tabular}{@{}lrr@{}}
\toprule
\textbf{Category} & \textbf{Time (ms)} & \textbf{Share (\%)} \\
\midrule
ZKP proving (Invoice + Transfer) & 464.6 & 46.0 \\
ZKP verification (Invoice + Transfer) & 303.1 & 30.0 \\
Setup (Invoice + Transfer) & 181.4 & 18.0 \\
Auxiliary (anchor, decrypt, committee) &  60.4 & 6.0 \\
\midrule
\textbf{Total} & \textbf{1009.6} & \textbf{100} \\
\bottomrule
\end{tabular}
\caption{Cost breakdown by functional category.}
\label{tab:breakdown}
\end{table}

The two ZKP proving steps dominate at 46\% of the total: Invoice Prove (305\,ms) accounts for three Paillier exponentiations in $\ZZ_{N^2}$ (two commitment exponentiations plus one EC scalar multiplication), while Transfer Prove (159\,ms) requires two (one Paillier, one EC). Verification is the second-largest cost at 30\%, with a similar structure but operating on the verification equations rather than generating commitments. Setup operations (key derivation, Paillier encryption of the refresh value and holder secret) contribute 18\%. The remaining 6\% consists of the anchor-point verification ($<\!10\,\mu$s, a single EC point comparison), Bob's Paillier decryption (30\,ms), and the committee's share update (30\,ms, also a single Paillier decryption).

\paragraph{Bottleneck analysis.} The dominant cost is modular exponentiation in $\ZZ_{N^2}$ ($|N| = 3072$, so 6144-bit integers). EC operations (256-bit secp256k1) are roughly $10^3\times$ faster and contribute negligibly. This Paillier-heavy cost profile is why circuit-based backends, which must express these exponentiations as field-arithmetic constraints, pay orders-of-magnitude overhead.

\paragraph{Proof size.} Each Invoice Consistency Proof $\pi_{\mathrm{inv}}$ consists of three commitments ($A_D \in \GG$, $A_B, A_C \in \ZZ_{N^2}$), one shared response $z \in \ZZ_q$, and two randomness responses $\tau_B, \tau_C \in \ZZ_N$, totaling $\approx\!2.3$\,KB at $|N| = 3072$. Each Transfer Consistency Proof $\pi_{\mathrm{xfer}}$ totals $\approx\!1.5$\,KB. The combined proof payload per transfer step is therefore $\approx\!3.8$\,KB, small enough for a single TCP segment. A channel with $m$ transfers accumulates $\approx\!3.8m$\,KB of off-chain state ($\approx\!3.7$\,MB at 1000 transfers).


\subsection{System-Level Comparison: \thunderbolt vs.\ Lightning}\label{sec:eval-lightning}

\autoref{tab:lightning} summarizes the structural differences between \thunderbolt and the Lightning Network along the dimensions that matter for real-world deployment.

\begin{table}[t]
\centering
\small
\begin{tabular}{@{}lll@{}}
\toprule
\textbf{Dimension} & \textbf{Lightning} & \textbf{\thunderbolt} \\
\midrule
Topology & Routed graph & Per-UTXO holder chain \\
Connectivity & Path-reachable & No path required \\
Routing & Onion-routed HTLCs & No routing \\
Liquidity & Directional / channel & Entire UTXO moves \\
Rebalancing & Required & Not applicable \\
Online assumption & Both endpoints & Asynchronous \\
On-chain / channel & 2 tx (open, close) & 2 tx (open, close) \\
On-chain / transfer & 0 (amortized) & 0 (amortized) \\
Holders / channel & 2 (fixed endpoints) & Unbounded chain \\
Watchtowers & Required for safety & Committee (always-on) \\
Trust model & Trustless + WT & $(t,n)$ VESS \\
\bottomrule
\end{tabular}
\caption{System-level comparison between Lightning and \thunderbolt.}
\label{tab:lightning}
\end{table}

\autoref{tab:lightning} summarizes the structural differences between \thunderbolt and the Lightning Network along the dimensions that matter for real-world deployment.

\paragraph{Connectivity and routing.} Lightning's most binding constraint is that every payment requires a path in the channel graph; a payee with no inbound liquidity is unreachable. \thunderbolt eliminates this constraint entirely: a transfer from $\mathsf{H}_s$ to $\mathsf{H}_{s+1}$ requires only that $\mathsf{H}_{s+1}$ publishes an invoice.

\paragraph{Liquidity.} A Lightning channel's capacity is directional: funds pushed one way cannot flow back without rebalancing via on-chain or loop-out operations. In \thunderbolt, the entire UTXO transfers at once, so there is no directional capacity and no rebalancing cost.

\paragraph{Asynchrony.} Lightning HTLCs require both sender and receiver to be online for the duration of the payment. \thunderbolt transfers are asynchronous: Bob posts his invoice once; Alice consumes it whenever she chooses. This makes \thunderbolt suitable for payroll batches, cold-wallet receipts, and low-power devices that wake up only periodically.

\paragraph{Latency comparison.} Measured Lightning payment latencies vary widely by route length and network conditions. Using c-lightning probe-plugin data as reference: a 4-hop round-trip takes $\approx\!2.56$\,s, and the median payment completion time is $\approx\!13.1$\,s (including path-finding and retries). \thunderbolt completes a transfer in $\approx\!1.01$\,s, yielding \textbf{2.5$\times$}--\textbf{13$\times$} improvements. Crucially, \thunderbolt's latency is \emph{deterministic}---it depends on local Paillier computation, not network hops or failed routes.

\paragraph{Trust and on-chain scope.} Lightning is trustless modulo watchtowers. \thunderbolt relies on the Committee; the threshold-committee instantiation (\autoref{sec:vess-committee}) reduces this to standard $(t,n)$-threshold honesty. Both systems amortize on-chain cost, but Lightning amortizes per \emph{channel pair} ($\Theta(n)$ on-chain operations for $n$ users) while \thunderbolt amortizes per \emph{UTXO} ($O(1)$).


\subsection{Sigma Protocol vs.\ General-Purpose SNARKs}\label{sec:eval-snark}

An alternative realization of the Invoice and Transfer relations would encode them as arithmetic circuits and feed them into a general-purpose SNARK backend (e.g., a zkVM such as Succinct's SP1~\cite{sp1}, RISC Zero~\cite{risczero}, or a circuit-based backend such as Groth16~\cite{groth16} or Plonk~\cite{gabizon2019plonk}). We argue---and partially measure---that this path is orders of magnitude more expensive than our direct Sigma-protocol instantiation.

\paragraph{Cost drivers.} The Invoice and Transfer relations involve two kinds of expensive operations:
\begin{itemize}
  \item \textbf{Paillier arithmetic:} modular exponentiation in $\ZZ_{N^2}$ for a 3072-bit $N$, which in a SNARK circuit requires bit-decomposition of 6144-bit integers and thousands of field operations per exponentiation.
  \item \textbf{Elliptic-curve scalar multiplication:} secp256k1 scalar multiplication, requiring 256 doublings and up to 256 additions; each EC operation in a SNARK circuit is a non-trivial sub-circuit over a non-native field.
\end{itemize}
In a general-purpose circuit backend, each Paillier ciphertext verification alone consumes on the order of $10^5$--$10^6$ constraints, and each EC operation adds $10^3$--$10^4$ more. The resulting circuit size is dominated by Paillier exponentiations, making this one of the worst-case workloads for generic SNARK backends.

\paragraph{Sigma-protocol cost.} Our direct Sigma-protocol construction requires no circuit. Each proof consists of a constant number of native group operations:
\begin{itemize}
  \item Invoice Consistency Proof: 3 commitments (1~EC, 2~Paillier), 1~Fiat--Shamir hash, 1~shared response, 2~Paillier randomness responses. Verification: 3~equations (1~EC, 2~Paillier).
  \item Transfer Consistency Proof: 2 commitments (1~EC, 1~Paillier), 1~hash, 1~response, 1~Paillier randomness response. Verification: 2~equations (1~EC, 1~Paillier).
\end{itemize}
Concretely, proving both proofs takes $\approx\!465$\,ms and verifying both takes $\approx\!303$\,ms, with no constraint system, no trusted setup, and no circuit-level tooling. \autoref{tab:snark} summarizes the comparison.

\paragraph{SP1 benchmark methodology.} To provide a concrete reference for the SNARK approach, we implemented the same transfer semantics (holder-secret update $x_{s+1} = x_s + r_s$, committee update $y_{s+1} = y_s - r_s$, invariant check) inside Succinct's SP1 zkVM~\cite{sp1}. Note that this SP1 program uses native \texttt{u64} arithmetic, \emph{not} the full 3072-bit Paillier operations from our real protocol. It therefore represents an \emph{optimistic lower bound} on what a zkVM would cost for the full relation: the real SNARK circuit would additionally need to express Paillier modular exponentiation over 6144-bit integers, which is the dominant cost in our native benchmarks.

\begin{table}[t]
\centering
\small
\begin{tabular}{@{}lrrl@{}}
\toprule
\textbf{Backend} & \textbf{Prove} & \textbf{Verify} & \textbf{Setup} \\
\midrule
\thunderbolt (Sigma + FS) & 465\,ms & 303\,ms & None \\
SP1 zkVM (u64 only) & $\gg\!465$\,ms & $\gg\!303$\,ms & None \\
Groth16 (projected) & $\sim$min & $\sim$3\,ms & Trusted \\
\bottomrule
\end{tabular}
\caption{Proving and verification cost for one complete transfer step (Invoice + Transfer). SP1 row uses simplified \texttt{u64} arithmetic (optimistic lower bound). Groth16 row is projected from published Paillier-in-circuit estimates.}
\label{tab:snark}
\end{table}

\paragraph{Circuit complexity estimate.} The Invoice proof alone would exceed $10^6$ R1CS constraints ($\approx\!2 \times 10^4$ for EC scalar multiplication + $\approx\!3 \times 10^5$--$10^6$ per Paillier exponentiation over 6144-bit integers). At published Groth16 proving rates of $10^4$--$10^5$ constraints/s, this yields tens of seconds to minutes---a $100$--$1000\times$ overhead over our 305\,ms native prove. The zkVM path adds a further constant from RISC-V emulation.

\paragraph{Discussion.} Groth16 proofs are smaller (288~bytes vs.\ our $\approx\!3.8$\,KB) but require a per-circuit trusted setup; STARK-based backends are transparent but produce 50--200\,KB proofs. Our relations are \emph{algebraic} (DL + Paillier equalities with a shared witness), precisely the regime where Sigma protocols dominate circuit-based alternatives. \thunderbolt trades a modestly larger proof for orders-of-magnitude lower proving cost and no trusted setup.
